\scanpage{69}
% archive sheet number: 70
que d'autre part les valeurs propres non nulles de $u$ et de sa restriction $v$ à $E_0$ sont les mêmes, la prop. 3 est par là démontrée pour les op. $u$ de rang fini.

Comme on sait que le système \deleted{tout} entier, propre à un op. compact $u$ dans un espace de Banach, dépend \deleted{\unsure{continûment}} contin\unsure{û}ment \deleted{\unsure{pris comme un}} (dans un sens évident à préciser) de $u$, il s'ensuit que la prop. 3 est aussi vraie pour toute limite d'opérateurs de rang fini, i.e. pour tous les op. compacts.

On conclut de la prop. 3 les résultats suivants, qui seront essentiels au N\textsuperscript{o} suivant : \deleted{(\unsure{lettre Hubert|cette thèse} \unsure{pour} \illegible{2}}

\uline{Corollaire 2} Soit $u$ un op. \inserted{compact} dans un espace de Hilbert complexe $E$, \inserted{alors pour tout entier $n\geq1$, on a} \deleted{soit $(\lambda_i)$ la suite de ses valeurs propres non nulles, \unsure{chaque} \illegible{2} \unsure{pour} tout \unsure{$u\in L(E)$}, \unsure{$\lambda_i(\Lambda^n)$} pour un entier $n$, soit \unsure{$(\rho_i)$} la suite décroissante des valeurs propres de l'adjoint. Alors pour tout $n$ \unsure{on a}}

\[
|\lambda_1\ldots\lambda_n|\leq\rho_1\ldots\rho_n\tag*{(23)}
\]

\margin{au \uline{N\textsuperscript{o} 1}}

Voir conventions de notation \deleted{avant la prop 3} \inserted{au N\textsuperscript{o} 1}

En effet, d'après la \deleted{prop. 3 le \unsure{terme} est la plus grande des valeurs propres de} \inserted{\unsure{formule ci-contre} (23)} le premier membre est ég. à \deleted{\unsure{$\lambda_1$}} $|\lambda_1(\overset{n}{\Lambda}u)|$, donc \unsure{maj.} par $\lVert\overset{n}{\Lambda}u\rVert$, \deleted{\unsure{d'o}} ce qui n'est autre que le deuxième mb de (23), en vertu de (22).

\deleted{$\lVert\overset{n}{\Lambda}u\rVert=\lVert((\overset{n}{\Lambda}u)^*(\overset{n}{\Lambda}u))^{1/2}\rVert=\deleted{\unsure{(\overset{n}{\Lambda}(u^*u))^{1/2}}}\lVert\overset{n}{\Lambda}h\rVert$, où $h=(u^*u)^{1/2}$ (prop 1), or, $\lVert\overset{n}{\Lambda}h\rVert$ est \inserted{$\overset{n}{\Lambda}h$ étant hermitien positif, \unsure{par suite}} la plus grande des valeurs propres de $\overset{n}{\Lambda}h$, donc à $\rho_1\ldots\rho_n$ d'après (21).}

\uline{Corollaire 3} Soient \deleted{\unsure{$u$ et $v$ des op. compacts dans un}} \inserted{$E,F,G$ trois espaces de Hilbert,}

% margin note written sideways
\margin{\uline{Corollaire 1} Soit $u$ \deleted{un op. compact de l'un deux} \inserted{opérateur compact} de \inserted{$E$ dans un autre $F$} espace de Hilbert \deleted{\unsure{complexe}}, on a pour tout entier $\geq1$ \inserted{si $E=F$}}

% margin note written sideways
\margin{$\text{(21)}\quad\lambda_1(u)\ldots\lambda_n(u)=\lambda_1(\overset{n}{\Lambda}u)$}

% margin note written sideways
\margin{et dans tous les cas, on a $\text{(22)}\quad\rho_1(u)\ldots\rho_n(u)=\lVert\overset{n}{\Lambda}u\rVert$}

% margin note written sideways
\margin{\deleted{\unsure{on pose}} pour les notations, voir \unsure{la rédaction} qui précède la prop. 3). En effet, la formule (21) est une \unsure{conséquence} \unsure{immédiate} de la prop. 3. Quant à la formule (22), on note que $\rho_i(u)^2=\lambda_i(u^*u)$, donc le carré du premier mb de (22) est en vertu de (21), égal à}

% slanted
\margin{$\lambda_1(\overset{n}{\Lambda}(u^*u))=\rho_1((\overset{n}{\Lambda}u)^*\overset{n}{\Lambda}u)$ (\unsure{prop. 2} Voir \deleted{\illegible{2}}) $=\lVert\overset{n}{\Lambda}u\rVert^2$ (\unsure{facteurs $\rho$} \unsure{$h=\, $} \illegible{2}}
